\chapter{Tabella comparativa delle notifiche}
\label{cap:scadenze}

Il professore ha avvertito che all'esame va identificata prima la norma applicabile e poi la
tempistica corretta, senza confondere le fasi delle diverse norme. La tabella riassume le
tempistiche viste finora.

\begin{center}
\small
\begin{tabularx}{\linewidth}{>{\raggedright\arraybackslash}p{1.7cm}
  >{\raggedright\arraybackslash}X >{\raggedright\arraybackslash}X
  >{\raggedright\arraybackslash}X >{\raggedright\arraybackslash}X}
\toprule
\textbf{Norma} & \textbf{Evento e destinatario} & \textbf{Fase 1} & \textbf{Fase 2} & \textbf{Fase 3} \\
\midrule
\textbf{GDPR} (artt.~33--34)
  & Data breach $\rightarrow$ autorità di controllo (Garante; capofila se transfrontaliero)
  & Notifica entro \textbf{72 h dalla consapevolezza} (anche preliminare)
  & Completamento: \textbf{nessun termine} (prassi $\approx$ 1 mese)
  & Interessati \textbf{senza ingiustificato ritardo}, se rischio elevato (art.~34) \\
\midrule
\textbf{NIS2} (art.~23)
  & Incidente significativo $\rightarrow$ CSIRT (Italia: CSIRT Italia, ACN)
  & Early warning \textbf{24 h dalla consapevolezza}
  & Notifica \textbf{72 h dalla consapevolezza}; relazione intermedia su richiesta
  & Final report \textbf{1 mese dalla notifica} \\
\midrule
\textbf{DORA}
  & Incidente ICT grave $\rightarrow$ autorità competente di settore
  & Notifica iniziale \textbf{4 h dalla classificazione} come grave, max \textbf{24 h dalla conoscenza}
  & Relazione intermedia \textbf{72 h dalla notifica iniziale}
  & Relazione finale \textbf{1 mese} \\
\midrule
\textbf{CRA} (art.~14(2))
  & Vulnerabilità attivamente sfruttata $\rightarrow$ CSIRT coordinatore + ENISA (Single Reporting Platform)
  & Early warning \textbf{24 h dalla conoscenza}
  & Notifica \textbf{72 h dalla conoscenza}
  & Final report \textbf{14 giorni dalla disponibilità della misura correttiva} \\
\midrule
\textbf{CRA} (art.~14(4))
  & Incidente grave sulla sicurezza del prodotto $\rightarrow$ come sopra
  & Early warning \textbf{24 h}
  & Notifica \textbf{72 h}
  & Final report \textbf{1 mese dalla notifica} \\
\midrule
\textbf{AI Act} (art.~73)
  & Incidente grave causato da un sistema di IA ad alto rischio $\rightarrow$ autorità di vigilanza del mercato
  & Notifica \textbf{15 giorni} dalla conoscenza (stabilito il nesso causale)
  & \textbf{2 giorni} se infrastruttura critica o violazione diffusa
  & \textbf{10 giorni} in caso di morte \\
\bottomrule
\end{tabularx}
\end{center}

\begin{esame}
\textbf{Trappole tipiche:}
\begin{itemize}
  \item le \textbf{72 ore di NIS2} decorrono dalla \emph{consapevolezza}; le \textbf{72 ore di DORA}
    decorrono dalla \emph{notifica iniziale};
  \item il \textbf{final report CRA sulle vulnerabilità} decorre dalla \emph{disponibilità del fix},
    non dalla consapevolezza;
  \item il \textbf{GDPR non ha final report} con scadenza;
  \item l'\textbf{early warning CRA} non è la \textbf{notifica preliminare GDPR}: sono obblighi
    distinti, verso autorità distinte;
  \item una notifica fatta a un'autorità \textbf{non adempie} gli obblighi verso un'altra: se lo stesso
    evento rientra in più norme, le notifiche corrono in parallelo;
  \item le scadenze \textbf{corrono anche nel weekend}.
\end{itemize}
\end{esame}

\begin{verifica}
\textbf{DORA, relazione finale}: gli appunti dell'anno precedente (\emph{main.pdf}) riportano in un
punto ``1 mese dalla relazione intermedia o dal suo ultimo aggiornamento'' e in un altro ``1 mese
dall'incidente''. Verificare sulle slide DORA di quest'anno.
\end{verifica}

\begin{integrazione}[DORA, relazione finale]
Le norme tecniche di attuazione sulla notifica degli incidenti (Reg.~delegato (UE) 2025/301) fissano la
relazione finale \textbf{entro un mese dall'ultima relazione intermedia aggiornata}. Le stesse norme
consentono ad alcune entità, se la scadenza cade nel fine settimana o in un festivo, di inviare entro
mezzogiorno del primo giorno lavorativo: per DORA la regola ``le scadenze corrono anche nel weekend'' va
quindi presa con cautela. Da confermare con le slide DORA.
\end{integrazione}
